\section*{Resumen}
\addcontentsline{toc}{section}{Resumen}

Este informe presenta el segundo entregable del Trabajo Parcial de Programación Concurrente y
Distribuida y, como pide el enunciado, conserva el primero con sus correcciones. El caso es el
mismo: K-means sobre los viajes del taxi amarillo de Nueva York \parencite{nyctlc} para
caracterizar patrones de movilidad urbana, en línea con el Objetivo de Desarrollo Sostenible~11.

En este entregable el algoritmo se implementó en Go, sin librerías de terceros, en dos versiones
que comparten el mismo contrato numérico: una secuencial y una concurrente con un \emph{worker pool}
persistente de goroutines, un canal que reparte bloques de viajes, acumuladores privados por bloque
y una barrera con \texttt{sync.WaitGroup} al final de cada iteración. La lógica de sincronización se
modeló en Promela y se verificó con Spin sobre todos los entrelazados de un dominio reducido: el
modelo correcto termina sin errores y un mutante con un acumulador compartido, sin exclusión mutua,
produce el contraejemplo esperado.

El speedup se midió con un protocolo fijado en código: \cifra{repeticiones}~repeticiones por
configuración en orden aleatorio, media recortada al \cifra{recorte} por extremo e intervalos de
confianza por bootstrap. Sobre los \cifra{n}~viajes del dataset, la versión concurrente alcanza
\cifra{speedup-p4}$\times$ con cuatro workers y \cifra{speedup-p8}$\times$ con ocho, con la misma
inercia final para cualquier cantidad de workers. El análisis muestra que el costo que limita la
escalabilidad es el de coordinar, que crece con los workers, y que el paralelismo deja de pagar por
debajo de unos \cifra{n-equilibrio}~viajes, el punto de equilibrio.

\textit{Palabras clave:} K-means, programación concurrente, Go, worker pool, Promela, Spin, speedup,
media recortada, NYC TLC.
